\scanpage{95}
% archive foliation: 94
savoir les ``invariants birationnels fondamentaux''
\inserted{relativement à $\mathbf F_p$,}
\deleted{de} envisagés dans les notes sur les \uline{Motifs}.

Il semblerait opportun d'étudier \deleted{\unsure{si}}
\inserted{le cas des coefficients $F$ généraux et de voir si,}
pour un morphisme $f:X\longrightarrow Y$
disons lisse \deleted{projectif}\inserted{propre}, \inserted{$F$ \unsure{lct|ct} tordu sur $X$,} la connaissance de
\deleted{$\unsure{H|\zeta}$}
$L(X,F)$ modulo $L(f^{-1}(Y'),F)$,
pour $Y'$ propre dans $Y$, permet de récupérer
les invariants birationnels fondamentaux
de $(X_K,\unsure{F\vert_{X_K}|F\vert_K})/K$
[$K$ corps des fonctions de
$Y$]. \uline{Dans quelle mesure faut-il}
\uline{faire les hypothèses $f$ lisse propre}, et $F$
\uline{ct tordu} pour obtenir des énoncés
de cette nature ??
\figure{\begin{tikzpicture}[scale=.45]
\draw (0,0)--(-.04,2.3)--(-.08,4.6) (.16,0)--(.11,2.3)--(.09,4.6);
\end{tikzpicture}}
